%%%%%%%%%%%%%%%%%%%%%%%%%%%%%%%%%%%%%%%%
%% MCM/ICM LaTeX Template (Enhanced)  %%
%% 2026 MCM/ICM                        %%
%% Full 3-level headings + real English content
%%%%%%%%%%%%%%%%%%%%%%%%%%%%%%%%%%%%%%%%
\documentclass[12pt]{article}

% ---------- Page & Font ----------
\usepackage{geometry}
\geometry{left=1in,right=0.75in,top=1in,bottom=1in}
\usepackage{newtxtext}
\usepackage{amsmath,amssymb,amsthm}
\usepackage{newtxmath} % must come after amsXXX

% ---------- Figures / Tables ----------
\usepackage{graphicx}
\usepackage{float}
\usepackage{booktabs}
\usepackage{tabularx}
\usepackage{array}
\usepackage{subcaption}

% ---------- Lists ----------
\usepackage{enumitem}

% ---------- Citations ----------
\usepackage[numbers]{natbib}

% ---------- Colors / Hyperlinks ----------
\usepackage{xcolor}
\definecolor{MyBlue}{RGB}{0,72,153}
\definecolor{MyRed}{RGB}{150,0,0}
\usepackage{hyperref}
\hypersetup{colorlinks=true, linkcolor=MyBlue, urlcolor=MyRed, citecolor=MyBlue}



% ---------- Header / Footer ----------
\usepackage{fancyhdr}
\pagestyle{fancy}
\lhead{Team \Team}
\rhead{Page \thepage}
\cfoot{}

%%%%%%%%%%%%%%%%%%%%%%%%%%%%%%%%%%%%%%%%
% Replace ABCDEF in the next line with your chosen problem
% and replace 1111111 with your Team Control Number
\newcommand{\Problem}{ABCDEF}
\newcommand{\Team}{1111111}
%%%%%%%%%%%%%%%%%%%%%%%%%%%%%%%%%%%%%%%%

% ---------- Theorem-like environments ----------
\newtheorem{theorem}{Theorem}
\newtheorem{lemma}[theorem]{Lemma}
\newtheorem{definition}{Definition}

% ---------- Handy macros ----------
\newcommand{\R}{\mathbb{R}}
\newcommand{\E}{\mathbb{E}}
\newcommand{\Var}{\mathrm{Var}}
\newcommand{\Prob}{\mathbb{P}}
\newcommand{\vect}[1]{\boldsymbol{#1}}
\newcommand{\norm}[1]{\left\lVert #1 \right\rVert}
\newcommand{\argmin}{\operatorname*{arg\,min}}
\newcommand{\argmax}{\operatorname*{arg\,max}}

%%%%%%%%%%%%%%%%%%%%%%%%%%%%%%%%%%%%%%%%
\begin{document}
\graphicspath{{.}{figures/}}  % put images in ./figures or current dir
\DeclareGraphicsExtensions{.pdf,.jpg,.jpeg,.png}

%%%%%%%%%%%%%%%%%%%%%%%%%%%%%%%%%%%%%%%%
%              SUMMARY SHEET
%%%%%%%%%%%%%%%%%%%%%%%%%%%%%%%%%%%%%%%%
\thispagestyle{empty}
\vspace*{-16ex}
\centerline{\begin{tabular}{*3{c}}
  \parbox[t]{0.3\linewidth}{\begin{center}\textbf{Problem Chosen}\\ \Large \textcolor{red}{\Problem}\end{center}}
  & \parbox[t]{0.3\linewidth}{\begin{center}\textbf{2026\\ MCM/ICM\\ Summary Sheet}\end{center}}
  & \parbox[t]{0.3\linewidth}{\begin{center}\textbf{Team Control Number}\\ \Large \textcolor{red}{\Team}\end{center}}\\
  \hline
\end{tabular}}


\begin{center}
{\Large\textbf{Put your own title HERE}}\\
\vspace{0.4em}
\textbf{Summary}
\end{center}

We develop a two-stage framework that links prediction with decision-making under practical constraints. In Stage I,
we construct a parsimonious predictive model $\hat{y}=f(\vect{x};\theta)$ that maps scenario features to the target quantity of interest.
We calibrate parameters using least squares / maximum likelihood and validate performance via a hold-out test and stress scenarios.
In Stage II, we design a constrained optimization that selects an actionable policy $\vect{d}$ to maximize expected utility while respecting
budget, feasibility, and fairness constraints.
Across our benchmark scenarios, the proposed policy improves the primary score by 12.4\% compared with a baseline heuristic,
while keeping the violation rate below 1\%. Sensitivity analysis shows that conclusions remain stable under $\pm 10\%$
perturbations of the most influential parameters.
We recommend deploying the optimized policy with a conservative safety margin and updating parameters periodically as new data arrive.

\vspace{0.8em}
\noindent\textbf{Key words:} prediction; optimization; robustness; validation

\clearpage

%%%%%%%%%%%%%%%%%%%%%%%%%%%%%%%%%%%%%%%%
%                MAIN BODY
%%%%%%%%%%%%%%%%%%%%%%%%%%%%%%%%%%%%%%%%
\setcounter{page}{1}
% Uncomment if you want a TOC (counts toward page limit)
\setcounter{tocdepth}{2}
\tableofcontents

\clearpage

%%%%%%%%%%%%%%%%%%%%%%%%%%%%%%%%%%%%%%%%
\section{Introduction}
\subsection{Background and Motivation}
Many real-world planning problems require both (i) forecasting uncertain outcomes and (ii) choosing actions that perform well
under constraints. A purely predictive model may be accurate yet operationally useless, while a purely heuristic decision rule
may be feasible yet suboptimal. Our goal is to bridge this gap by combining a lightweight predictive model with an optimization layer
that explicitly encodes feasibility and risk preferences.

\subsubsection{Quick overview of our approach}
Our approach follows the pipeline:
\[
\text{data/assumptions} \rightarrow \text{feature design} \rightarrow \text{prediction} \rightarrow \text{optimization} \rightarrow \text{validation/sensitivity}.
\]
For official contest information, see \href{https://www.contest.comap.com/undergraduate/contests/mcm/}{COMAP MCM/ICM website}.

\section{Problem Restatement}
We restate the problem in terms of inputs, outputs, and evaluation criteria:
\begin{itemize}[leftmargin=*]
  \item \textbf{Inputs:} observed historical records, scenario descriptors, and constraints (budget/capacity/feasibility).
  \item \textbf{Outputs:} forecasts for target quantities and a recommended policy/strategy.
  \item \textbf{Objectives:} maximize performance score (or minimize cost/risk) subject to stated constraints.
\end{itemize}

\subsubsection{Task decomposition}
We decompose the problem into three tasks:
\begin{enumerate}[leftmargin=*]
  \item \textbf{Task 1 (Predict):} estimate outcomes $\hat{y}$ from features $\vect{x}$.
  \item \textbf{Task 2 (Decide):} choose an action/policy $\vect{d}$ that optimizes an objective under constraints.
  \item \textbf{Task 3 (Explain):} interpret results, provide insights, and assess robustness.
\end{enumerate}

\section{Assumptions}
Assumptions are designed to be defensible and to simplify modeling without breaking realism:
\begin{enumerate}[leftmargin=*]
  \item \textbf{Stable mapping:} the relationship between features and outcomes is approximately stable over the planning horizon.
  \item \textbf{Bounded noise:} measurement errors are bounded and do not systematically bias the main outcomes.
  \item \textbf{Constraint validity:} budgets/capacities provided in the prompt reflect operational limits.
\end{enumerate}

\subsubsection{What we do if an assumption fails}
If stability is violated (distribution shift), we re-fit using rolling windows; if noise is heavy-tailed, we switch from least squares
to a robust loss (e.g., Huber). These contingency plans are tested via stress scenarios in Section~\ref{sec:sensitivity}.

%%%%%%%%%%%%%%%%%%%%%%%%%%%%%%%%%%%%%%%%
\section{Notation and Definitions}
\subsection{Notation table}
\begin{table}[H]
\centering
\caption{Key notation used throughout the paper.}
\label{tab:notation}
\begin{tabularx}{0.95\linewidth}{>{\centering\arraybackslash}m{0.20\linewidth} X}
\toprule
\textbf{Symbol} & \textbf{Meaning} \\
\midrule
$\vect{x}\in\R^n$ & Feature vector describing a scenario or candidate decision \\
$y\in\R$ & Observed target value (score/cost/demand/etc.) \\
$\hat{y}=f(\vect{x};\theta)$ & Prediction model with parameters $\theta$ \\
$\vect{d}\in\R^m$ & Decision/policy variables to be selected \\
$A\vect{d}\le \vect{b}$ & Linear feasibility constraints (capacity/budget) \\
$\lambda,\gamma$ & Trade-off parameters for utility vs. cost/risk \\
\bottomrule
\end{tabularx}
\end{table}

\subsection{Core definition and a lemma}
\begin{definition}[Feasible set]
The feasible set is $\mathcal{F}=\{\vect{d}\in\R^m: A\vect{d}\le \vect{b},\ \vect{d}\ge \vect{0}\}$.
\end{definition}

\begin{lemma}[Convex feasibility]
If $A\vect{d}\le\vect{b}$ and $\vect{d}\ge\vect{0}$ are linear constraints, then $\mathcal{F}$ is convex.
\end{lemma}

%%%%%%%%%%%%%%%%%%%%%%%%%%%%%%%%%%%%%%%%
\section{Overall Solution Framework}
\subsection{Workflow diagram}
Figure~\ref{fig:workflow} illustrates the end-to-end pipeline from data to decisions and evaluation.

\begin{figure}[H]
  \centering
  % Replace with your own figure file, e.g., figures/workflow.png
  \includegraphics[width=0.85\linewidth]{example-image-a}
  \caption{Workflow: preprocessing $\rightarrow$ prediction $\rightarrow$ optimization $\rightarrow$ validation \& sensitivity.}
  \label{fig:workflow}
\end{figure}

\subsection{What makes the framework contest-friendly}
We prioritize (i) interpretability (clear assumptions), (ii) computational simplicity (fast fitting and solving),
and (iii) robustness (sensitivity tests). This balance helps produce a solution that is not only numerically strong
but also easy for judges to audit.

\subsubsection{Deliverables}
Our final deliverables include:
\begin{itemize}[leftmargin=*]
  \item a predictive model with reported error metrics,
  \item an optimization-based policy with constraint guarantees,
  \item a robustness analysis summarizing where conclusions may change.
\end{itemize}

%%%%%%%%%%%%%%%%%%%%%%%%%%%%%%%%%%%%%%%%
\section{Task 1: Predictive Modeling}
\subsection{Data preprocessing}
\subsubsection{Cleaning and normalization}
We remove duplicates, handle missing values by domain-appropriate imputation (median for continuous variables, mode for categorical),
and normalize continuous features to zero mean and unit variance so coefficients are comparable.

\subsubsection{Feature engineering}
We derive stable, interpretable features such as relative change:
\begin{equation}
\Delta x_t=\frac{x_t-x_{t-1}}{\max\{x_{t-1},\epsilon\}},\quad \epsilon>0.
\label{eq:relchange}
\end{equation}
This reduces sensitivity to scale and improves generalization.

\subsection{Model specification}
\subsubsection{Baseline regression}
We start with a parsimonious linear model:
\begin{equation}
\hat{y}=\beta_0+\sum_{j=1}^{n}\beta_j x_j.
\label{eq:lin}
\end{equation}

\subsubsection{Parameter estimation}
Parameters are obtained by least squares:
\begin{equation}
\hat{\beta}=\argmin_{\beta}\sum_{i=1}^{N}\big(y_i-\hat{y}_i(\beta)\big)^2.
\label{eq:ls}
\end{equation}

\subsection{Model evaluation}
\subsubsection{Train/validation protocol}
We split data into training and validation sets (or use $K$-fold cross-validation) to estimate out-of-sample performance.

\subsubsection{Metrics}
We report MAE and MAPE:
\[
\mathrm{MAE}=\frac{1}{N}\sum_{i=1}^{N}|y_i-\hat{y}_i|,
\qquad
\mathrm{MAPE}=\frac{100}{N}\sum_{i=1}^{N}\left|\frac{y_i-\hat{y}_i}{y_i}\right|.
\]

%%%%%%%%%%%%%%%%%%%%%%%%%%%%%%%%%%%%%%%%
\section{Task 2: Decision Model and Optimization}
\subsection{Decision variables and objective}
\subsubsection{Utility formulation}
We select a decision $\vect{d}$ to maximize expected utility minus cost:
\begin{equation}
\max_{\vect{d}\in\mathcal{F}}\; U(\vect{d})-\gamma\,C(\vect{d}),
\label{eq:utility}
\end{equation}
where $U(\vect{d})=\E[\hat{y}(\vect{d})]$ and $C(\vect{d})$ is a measurable cost.

\subsubsection{Concrete example objective}
A common contest-friendly choice is a linearized objective:
\[
U(\vect{d})=\vect{w}^\top \vect{d},\qquad C(\vect{d})=\vect{c}^\top \vect{d}.
\]
Then Eq.~\eqref{eq:utility} becomes a linear program.

\subsection{Constraints}
\subsubsection{Budget and capacity}
We enforce:
\begin{equation}
A\vect{d}\le\vect{b},\qquad \vect{d}\ge \vect{0}.
\label{eq:constraints}
\end{equation}

\subsubsection{Fairness / balance (optional)}
If needed, we add constraints such as $d_i\le \rho \sum_j d_j$ to prevent over-concentration.

\subsection{Solution method}
\subsubsection{Solver choice}
Because the optimization is linear/convex, we can solve it efficiently using standard LP/QP solvers.
This is computationally stable and reproducible.

\subsubsection{Sanity checks}
We verify feasibility (no violated constraints) and compare with baseline heuristics (e.g., greedy or uniform allocation).

%%%%%%%%%%%%%%%%%%%%%%%%%%%%%%%%%%%%%%%%
\section{Results}
\subsection{Main quantitative results}
\subsubsection{Result table}
\begin{table}[H]
\centering
\caption{Illustrative results table (replace with your real numbers).}
\label{tab:results}
\begin{tabular}{lcc}
\toprule
Method & Primary Score $\uparrow$ & Constraint Violation $\downarrow$ \\
\midrule
Baseline heuristic & 0.71 & 2.5\% \\
Proposed framework & 0.80 & 0.8\% \\
\bottomrule
\end{tabular}
\end{table}

\subsubsection{Interpretation}
The proposed approach improves the primary score while reducing violations, indicating that the optimization layer effectively
translates predictive estimates into feasible decisions.

\subsection{Visual results}
\subsubsection{Comparison figure}
\begin{figure}[H]
  \centering
  % Replace with your own plot, e.g., figures/compare.png
  \includegraphics[width=0.75\linewidth]{example-image-b}
  \caption{Illustrative comparison plot (placeholder).}
  \label{fig:compare}
\end{figure}

\subsubsection{What the figure shows}
Figure~\ref{fig:compare} summarizes relative performance under multiple scenarios and highlights stability of the proposed policy.

%%%%%%%%%%%%%%%%%%%%%%%%%%%%%%%%%%%%%%%%
\section{Sensitivity Analysis}
\label{sec:sensitivity}
\subsection{Sensitivity design}
\subsubsection{One-at-a-time perturbation}
We perturb a key parameter (e.g., $\gamma$ or major coefficients in $\theta$) by $\pm 10\%$ and recompute outcomes.

\subsubsection{Scenario stress tests}
We generate extreme scenarios (high demand, low capacity, noisy measurements) to test failure modes.

\subsection{Sensitivity results}
\subsubsection{Stability statement}
Performance remains within a narrow band for moderate perturbations, suggesting the decision is not fragile.

\subsubsection{Most influential factors}
We identify which parameter shifts cause the largest change and recommend monitoring those in deployment.

%%%%%%%%%%%%%%%%%%%%%%%%%%%%%%%%%%%%%%%%
\section{Strengths and Weaknesses}
\subsection{Strengths}
\subsubsection{Interpretability}
The model uses transparent assumptions, a clear notation table, and objective functions aligned with the task.

\subsubsection{Reproducibility}
Every step (preprocessing, fitting, optimization, evaluation) is deterministic given inputs, making the solution easy to audit.

\subsection{Weaknesses}
\subsubsection{Data limitations}
If the available data are sparse or biased, predictive accuracy may degrade in unseen regimes.

\subsubsection{Model mismatch}
A simple model may miss nonlinear effects; we mitigate this by sensitivity tests and by suggesting a robust loss if needed.

%%%%%%%%%%%%%%%%%%%%%%%%%%%%%%%%%%%%%%%%
\section{Conclusion}
\subsection{Summary of what we built}
\subsubsection{Task-level recap}
Task 1 provides a validated predictor; Task 2 converts predictions into a feasible policy; Task 3 explains trade-offs and robustness.

\subsubsection{Actionable recommendation}
We recommend deploying the optimized policy with a conservative safety margin and re-estimating parameters periodically.

\subsection{Future work}
\subsubsection{Model upgrades}
Future extensions include nonlinear predictors, Bayesian uncertainty quantification, and adaptive re-optimization under streaming data.

\subsubsection{Operational improvements}
If real-time feedback is available, we can implement rolling updates and early-warning thresholds.

%%%%%%%%%%%%%%%%%%%%%%%%%%%%%%%%%%%%%%%%
\clearpage
\bibliographystyle{plainnat}
\begin{thebibliography}{99}
\bibitem[Author(Year)]{key}
Author. (Year). \textit{Title}. Journal/Conference, volume(issue), pages.

\bibitem[Hastie et~al.(2009)]{hastie2009}
Hastie, T., Tibshirani, R., and Friedman, J. (2009).
\textit{The Elements of Statistical Learning}. Springer.
\end{thebibliography}

%%%%%%%%%%%%%%%%%%%%%%%%%%%%%%%%%%%%%%%%
\appendix
\section{Appendix A: Additional Derivations}
\subsection{Derivation details}
\subsubsection{Example bound}
Assume $\norm{\nabla f(\vect{x})}\le L$. Then for any $\vect{x},\vect{z}\in\R^n$,
\[
|f(\vect{x})-f(\vect{z})|\le L\norm{\vect{x}-\vect{z}}.
\]
This supports the stability discussion in Section~\ref{sec:sensitivity}.
\clearpage
\section*{Report on Use of AI Tools}
\addcontentsline{toc}{section}{Report on Use of AI Tools}

\subsection*{Tools Used}
We used the following AI tool(s) during the preparation of this report:
\begin{itemize}
  \item \textbf{Model/Tool:} \textcolor{red}{(e.g., ChatGPT / GPT-4 / other)}.
  \item \textbf{Access Mode:} \textcolor{red}{(e.g., web interface / local deployment)}.
\end{itemize}

\subsection*{Purpose and Scope of Use}
The AI tool(s) were used for:
\begin{itemize}
  \item \textbf{Language refinement:} improving grammar, clarity, and conciseness of sentences written by the team.
  \item \textbf{Formatting assistance:} generating LaTeX snippets (tables, figures, references) without changing the mathematical meaning.
  \item \textbf{Code sanity-checking:} spotting potential bugs or edge cases in team-written code (final code authored and executed by the team).
\end{itemize}

\subsection*{Human Verification and Responsibility}
All AI-generated outputs were treated as suggestions. The team verified and finalized all content by:
\begin{itemize}
  \item cross-checking every numerical result with independent calculations and/or rerunning code,
  \item confirming that assumptions match the problem statement and do not introduce prohibited information,
  \item ensuring that all claims supported by external knowledge are cited to non-AI sources.
\end{itemize}
The team takes full responsibility for the correctness, originality, and integrity of the final submission.

\subsection*{Record of AI-Generated Contributions}
Examples of AI-assisted elements include:
\begin{itemize}
  \item draft alternative phrasings of the Summary section,
  \item proposed outline structures for sections and subsections,
  \item LaTeX formatting suggestions for tables/figures and hyperlink settings.
\end{itemize}

\subsection*{Sources and Citations}
When AI outputs referenced facts or methods beyond the team's direct derivations, we validated them against reputable sources and cited those sources in the bibliography.


\end{document}
